\section{Adjunction on the entire dlt floor}\label{sec:adjunction}

We now prove the boundary-generation statement applied to the model
supplied by Proposition~\ref{prop:supported-model}.  The adjoint line already exists
on the reduced floor as a restriction from the ambient space.  Our task is
to construct enough sections of that line which agree through all of its
intersections.  This section sets up the strata and the actual residue
identities; the next three sections construct the compatible sections.

\begin{theorem}[Generation on the whole dlt floor]\label{thm:floor-generation}
Let \((V,B)\) be a normal ordinary \(\Q\)-factorial compact K\"ahler dlt
fourfold with effective rational boundary and analytically nef actual
adjoint \(A=K_V+B\).  Suppose that it has a projective log resolution
\(\pi:\widehat V\to V\) for which the strict boundary and exceptional
support have globally smooth distinct SNC components, every exceptional
crepant coefficient is below one, and \(\pi\) is generically an
isomorphism on the image of every irreducible component of an intersection
of distinct strict transforms of coefficient-one boundary primes.
Then the actual restriction \(A|_S\) to the entire reduced floor
\(S=\floor B\) is semiample.
\end{theorem}

If \(S\) is empty the conclusion is vacuous.  Otherwise its components
are three-dimensional.  We may work on each connected component of \(V\)
and take a common multiple over the finitely many components, so in the
proof we assume \(V\) connected and hence irreducible.  The theorem needs no section of the ambient
adjoint.  Its resolution hypothesis is precisely the output of
Lemma~\ref{lem:special-resolution}, so it does not add a hypothesis to
Theorem~\ref{thm:main}.

\subsection{Local adjunction calculations}

Lemma~\ref{lem:floor-connectedness} supplies the connectedness and
rationality properties used for successive strata.  We next compare an
adjunction coefficient on a normal surface slice.

\begin{lemma}[A normal surface slice for an adjunction coefficient]
\label{lem:adjunction-slice}
Let \(Z\) be a normal Cohen--Macaulay irreducible analytic space of
dimension \(n\geq2\), let \(D\) be a normal prime divisor, and let
\(B=D+B'\) be an effective rational boundary with actual
\(\Q\)-Cartier adjoint.  Fix a divisible degree \(m\), a local frame of
\(\OO_Z(m(K_Z+B))\), and its meromorphic \(m\)-residue on \(D\).
For a prime \(\Gamma\subset D\), a general point of \(\Gamma\) admits a
local slice by \(n-2\) holomorphic parameters with the following properties.
The slice \(T\) is a normal surface germ, the cut \(C=D\cap T\) is a
smooth curve germ transverse to \(\Gamma\), and the restricted line is
the actual \(m\)-adjoint line of the effective sliced boundary on \(T\).
After division by the base parameter volume, the order along \(\Gamma\)
of the residue on \(D\) equals the order at \(C\cap\Gamma\) of the
surface residue on \(C\).  For \(n=2\), the slice is the surface germ
itself.
\end{lemma}
\begin{proof}
Work in a relatively compact local embedding in \(\C^N\), and refine
the regular loci of \(Z\), its singular locus, \(D\), \(\Gamma\), the
boundary primes, and their intersections into a locally finite collection
of smooth strata.  After a neighborhood shrink only finitely many relevant
strata meet the chosen compact closure.  Choose a linear map
\(\ell:\C^N\to\C^{n-2}\) whose differential is an isomorphism on
the tangent space of \(\Gamma\) at a general smooth point where \(D\)
is smooth.  Such points exist because \(D\) is normal.  For the restriction
of \(\ell\) to each smooth stratum of dimension at least \(n-2\),
Sard's theorem makes its critical values a measure-zero subset of the
parameter space; strata of smaller dimension have measure-zero image
\cite{Sard1942}.  The image of \(\Gamma\) contains a neighborhood of the
chosen value.  We may therefore choose a nearby value regular for all
these restrictions and still meeting \(\Gamma\) at a general point.
The corresponding level is transverse to every stratum it meets.

At regular points of \(Z\) the level is a smooth surface.  The singular
locus of a normal \(n\)-fold has dimension at most \(n-2\); its strata
therefore meet this level only discretely.  Every irreducible component
of the level has dimension at least two by the principal ideal theorem.
It cannot be contained in that discrete singular intersection, and hence
has dimension exactly two by its smooth dense part.  The \(n-2\) level
equations have height \(n-2\) at every point.  In the Cohen--Macaulay local
rings of \(Z\) they are a regular sequence.  The quotient is consequently
Cohen--Macaulay of dimension two.  It is generically reduced and regular
in codimension one, since its possible singular points are discrete;
Serre's criterion makes it normal.  Transversality on \(D\) gives a smooth
curve at the cut point, and the fact that \(d\ell\) is an isomorphism on
\(T\Gamma\) makes this curve transverse to \(\Gamma\).

Each boundary prime cuts an effective curve with its original rational
coefficient, counted with its positive intersection multiplicity.  At a
generic codimension-one point of the normal surface, the ambient space
and the slice are smooth and transverse to these primes.  Dividing the
ambient canonical form by the parameter volume identifies the
restriction of the frame with the pluriadjoint frame of this sliced
effective boundary.  Both are rank-one reflexive sheaves on the normal
surface, so the identification extends from codimension one and is an
actual line identity.

To compare the orders, on the smooth \(D\) near a general point of
\(\Gamma\) complete the parameters to coordinates \((z_1,\ldots,z_{n-2},u)\)
with \(\Gamma=(u=0)\).  Write the meromorphic residue as
\[
 u^c h(z,u)(du\wedge dz_1\wedge\cdots\wedge dz_{n-2})^{\otimes m},
\]
where \(h\) is a unit at a general point of \(\Gamma\).  Choose the
regular level above also outside the proper zero set of its leading
coefficient.  Division by the base volume and restriction to that level
then has order exactly \(c\) on the cut curve.  On the ambient smooth
locus the two operations commute: in coordinates with \(D=(x=0)\),
both take the residue of
\(h(x,u,z)(dx/x\wedge du\wedge dz)^{\otimes m}\) to
\(h(0,u,z)(du)^{\otimes m}\) on the level, up to the ordering sign.
The preceding reflexive identification makes the surface term the residue
of the same actual frame.  The restricted original residue is already a
meromorphic form on the smooth curve \(C\), so equality on its dense
smooth-ambient part extends across the cut point.  This proves the order
assertion.  In even degree the ordering sign is one.
\end{proof}

\subsection{The indexed strata and their actual adjoints}

Fix the resolution in Theorem~\ref{thm:floor-generation}, and let
\(\widehat S_i\), for \(i\) in a finite set \(I\), be the strict transforms
of the components of \(S\).  For a nonempty subset \(J\subset I\), each
irreducible component \(\widehat Z\) of \(\bigcap_{i\in J}\widehat S_i\)
is smooth of dimension \(4-|J|\).  Call its reduced image
\(Z=\pi(\widehat Z)\) a stratum and write \(\pi_Z:\widehat Z\to Z\).
We also include \(V\), with the empty index set and resolution \(\widehat V\).
The special resolution makes each \(\pi_Z\) bimeromorphic and generically
an isomorphism.  The image determines its index set and its resolving
component: on the generic isomorphism locus, two different choices would
force an SNC intersection to have two different codimensions or local
branches.  There are finitely many strata, since the intersections are
compact analytic spaces.  Imposing one additional index \(j\notin J\)
cuts \(\widehat Z\) in a smooth, possibly disconnected divisor.  Its
components give the incidences from \(Z\) to the next strata.

To perform adjunction inductively, allow each unused floor coefficient
either to remain one or to become \(1/2\).  More precisely, for a stratum
indexed by \(J\), choose \(e_i=1\) for \(i\in J\) and
\(e_i\in\{1,1/2\}\) for \(i\notin J\), and put
\[
 B^{\mathbf e}=B-\sum_{i\in I}(1-e_i)S_i.
\]
Ordinary global \(\Q\)-factoriality makes this an actual rational-line
operation.  If \(\widehat G\) is the crepant boundary of \(B\), the new
one is
\[
 \widehat G^{\mathbf e}=\widehat G-
 \pi^*\sum_{i\in I}(1-e_i)S_i.
\]
The subtracted pullback is effective.  Exceptional coefficients remain
below one, and the coefficient-one primes are exactly the retained strict
floor primes.  Its support stays inside the resolved SNC support.

SNC adjunction along the indices in \(J\) defines a subboundary
\(\widehat B_{\widehat Z}^{\mathbf e}\) on \(\widehat Z\) and an actual
residue identity
\begin{equation}\label{eq:snc-stratum-adjunction}
 (K_{\widehat V}+\widehat G^{\mathbf e})|_{\widehat Z}
 =K_{\widehat Z}+\widehat B_{\widehat Z}^{\mathbf e}.
\end{equation}
Every coefficient is at most one.  Its coefficient-one divisors are
precisely the one-index intersections for unused indices with \(e_i=1\).
The equality means equality of the meromorphic maps from the same
restricted pluriadjoint line in a sufficiently divisible even degree.

\begin{proposition}[Adjunction and full-floor descent]\label{prop:strata-adjunction}
Every stratum \(Z\) is normal and compact K\"ahler.  For every allowed
choice \(\mathbf e\) there is an effective rational boundary
\(B_Z^{\mathbf e}\) such that
\begin{align}
 \pi_Z^*(K_Z+B_Z^{\mathbf e})
   &=K_{\widehat Z}+\widehat B_{\widehat Z}^{\mathbf e},
       \label{eq:stratum-crepant}\\
 \OO_Z\bigl(m(K_Z+B_Z^{\mathbf e})\bigr)
   &\simeq\OO_V\bigl(m(K_V+B^{\mathbf e})\bigr)|_Z
       \label{eq:stratum-line}
\end{align}
for all sufficiently divisible even \(m\).  Both identities are the
actual meromorphic residue identifications.  For the full weights write
\(B_Z=B_Z^{\mathbf 1}\) and \(C_Z=\floor{B_Z}\).  The pair \((Z,B_Z)\)
is dlt; its floor components are exactly the next incident strata.
In a common degree, sections on those components that agree by residue
on every subordinate stratum descend uniquely to a section on the entire
reduced \(C_Z\).  The same assertion holds for the floor \(S\subset V\).
\end{proposition}
\begin{proof}
We induct on the number of indices.  The assertions for the empty index
set are the given normality, effective boundary, and crepant identity on
\(V\).  Assume them for a stratum \(Z\) and all allowed weights there.
Lower all unused floor weights.  The resulting crepant SNC boundary on
\(\widehat Z\) has every coefficient below one.  The effective pair
\((Z,B_Z^{\mathbf e})\) is therefore klt.  Lemma~\ref{lem:floor-connectedness}
shows that \(Z\) has rational singularities and is Cohen--Macaulay.

To add an index \(j\), retain only that unused weight at one and lower
the others.  The coefficient-one locus on \(\widehat Z\) is
\(R=\widehat Z\cap\widehat S_j\), a disjoint union of smooth components.
Apply Lemma~\ref{lem:floor-connectedness} to \(\pi_Z\).  Its connected
fibers prevent the images of two such components from meeting.  On each
image \(D\), the equality \(\pi_{Z,*}\OO_R=\OO_{\pi_Z(R),\red}\)
identifies the proper bimeromorphic pushforward of the structure sheaf
of its smooth resolving component \(\widehat D\) with \(\OO_D\).
Factoring through the finite normalization shows that \(D\) is normal.
It is a compact analytic subspace of \(V\), so it inherits a K\"ahler
form by restricting local strictly plurisubharmonic potentials.

For arbitrary allowed weights retaining \(j\), define
\(B_D^{\mathbf e}=\pi_{D,*}\widehat B_{\widehat D}^{\mathbf e}\).
At this point it is a rational divisor whose effectivity remains to be
shown.  In a divisible even degree, a local frame of the ambient adjoint
line pulls back and takes SNC residue to a meromorphic pluriform on
\(\widehat D\).  The proper bimeromorphic map \(\widehat D\to D\) is an
isomorphism at the generic points of divisors of the normal \(D\).
There the same map identifies the divisorial adjoint for
\(B_D^{\mathbf e}\) with the restriction of the existing invertible line.
Normal reflexive extension gives \eqref{eq:stratum-line} on all of \(D\).
Pulling the line back to \(\widehat D\) recovers the original meromorphic
map, since the maps agree on a dense open.  This proves
\eqref{eq:stratum-crepant}, including its meromorphic interpretation.

We verify that \(B_D^{\mathbf e}\) is effective.  It is enough to test
the coefficient at a general point of each prime in \(D\), a codimension-two
locus in \(Z\).  Apply Lemma~\ref{lem:adjunction-slice}, using the
Cohen--Macaulay property already proved for \(Z\).  It reduces exactly
that coefficient, with the actual residue order preserved, to the smooth
cut curve in a normal surface germ with effective boundary.  On a minimal
resolution of this surface germ, write its
crepant boundary as the effective strict transform plus an exceptional
divisor \(E\).  For every exceptional curve \(C\),
\[
 E\cdot C=-(K_{\widetilde T}+B^{\mathrm{str}})\cdot C\leq0.
\]
Indeed \(B^{\mathrm{str}}\cdot C\geq0\), while surface adjunction gives
\(K_{\widetilde T}\cdot C=2p_a(C)-2-C^2\geq0\): exceptional
self-intersections are negative and a smooth rational exceptional
\((-1)\)-curve is excluded by minimality.  The negative-definite
exceptional intersection matrix, with nonnegative off-diagonal entries,
then implies \(E\geq0\).  Explicitly, if \(E=P-N\) has disjoint
nonnegative parts and \(N\ne0\), then
\(E\cdot N=P\cdot N-N^2>0\), contradicting \(E\cdot C\leq0\) on
each component of \(N\).  The strict
transform of the smooth cut curve is finite birational over it and hence
isomorphic.  Adjunction to that smooth curve in a smooth surface with
effective remaining boundary has a nonnegative coefficient.  This is
the coefficient originally tested.  There is no coefficient to test when
the new stratum is a point.  Effectivity follows and completes this part
of the induction.

For clarity, the full-weight pair on each stratum is dlt and has exactly
the claimed floor.  The SNC subboundary has coefficients at most one, so
it is lc, and crepancy makes the pair on \(Z\) lc.  On the smooth
\(\widehat Z\), let a divisorial
valuation have center of codimension \(c\), and choose local coordinates
\(x_1,\ldots,x_c,\ldots\) at a general point of that center, with the
SNC boundary components among the coordinate divisors.  Give an unmarked
coordinate weight zero and write the boundary weights as \(b_i\leq1\).
The reduced coordinate arrangement is lc, so
\[
 a(v;\widehat Z,\widehat B_{\widehat Z})
 \geq\sum_{i=1}^c(1-b_i)v(x_i).
\]
All \(v(x_i)\) here are positive.  A zero discrepancy therefore forces
every \(b_i=1\); its center is an intersection of coefficient-one
components.  Each such intersection has image meeting the isomorphism
locus of the original special resolution.  In particular each
coefficient-one divisor is nonexceptional for \(\pi_Z\); the others have
coefficients below one.  Let \(F_Z\subset Z\) be the complement of the
largest open over which \(\pi_Z\) is an isomorphism.  It has codimension
at least two on the normal target \(Z\), and contains no entire image of
a zero-discrepancy center.  Its inverse image on \(\widehat Z\) may already
be divisorial.  If necessary, principalize the pulled reduced ideal
\(\mathscr I(F_Z)\OO_{\widehat Z}\), and then apply the ordered-boundary
resolution as in Lemma~\ref{lem:special-resolution}, marking each distinct
prime in the prior SNC support, the newly principalized support, and every
divisorial exceptional prime.  All new exceptional centers lie over
\(F_Z\).  The displayed inequality makes their discrepancies positive.
This is a dlt resolution in the
standard sense.  It also proves that the coefficient-one primes on \(Z\)
are exactly the images of the one-index intersections.

It remains to check the asserted descent of sections.  With full weights,
let \(R=\widehat B_{\widehat Z}^{=1}\).  The preceding floor identification
gives \((\pi_Z(R))_{\red}=C_Z\); write \(\pi_R:R\to C_Z\) for the induced
restriction of \(\pi_Z\).  Sections on the incident strata
pull back to sections on its smooth components.  Residue compatibility
means equality on the entire scheme-theoretic pairwise intersections,
which are reduced smooth SNC strata.  The local equalizer for an SNC union
therefore glues them to a section on \(R\); it is the elementary equalizer
for the coordinate ideals of its components.  Even-degree iterated
residues make all paths to deeper intersections identical.  By
Lemma~\ref{lem:floor-connectedness} and projection formula,
\[
 (\pi_R)_*\pi_R^*(L|_{C_Z})=(L|_{C_Z})
\]
for any ambient Cartier line \(L\) under consideration.  The section
descends uniquely to all of \(C_Z\).  This reasoning applies also to
\(Z=V\) and \(C_Z=S\).  In particular it retains the conductor
coefficient at a general double crossing on each normalized component;
it does not posit an additional conductor boundary on the nonnormal
union itself.
\end{proof}

The proposition supplies a finite collection of normal compact K\"ahler
dlt strata, all carrying restrictions of the same actual ambient line.
It also specifies exactly what compatibility is needed to descend to a
whole floor.  The next section proves that, after imposing a small list of
birational residue comparisons on the lower strata, these compatible
boundary sections extend to their parent stratum.
